\subsection{Self-Adjoint Operators and Orthogonal Projections}

\begin{definition} \label{def:self-adjoint-hermitian} \hlt{Self-Adjoint and Hermitian}\\
An operator $T\in\LL(\mathcal{V})$ on a finite-dimensional inner-product space is \hlt{self-adjoint} if $T^*=T$. A square matrix $A$ is \hlt{Hermitian} if $A^*=A$. In orthonormal coordinates, these are the same condition by Theorem \ref{thm:matrix-of-adjoint-is-conjugate-transpose}.
\end{definition}

\begin{lemma} \label{lem:complex-zero-quadratic-form-implies-zero-operator} \hlt{Zero Complex Quadratic Form}\\
Let $\mathcal{V}$ be a finite-dimensional complex inner-product space and let $T\in\LL(\mathcal{V})$.\\
If $\innerprod{T\mathbf{v},\mathbf{v}}=0$ for every $\mathbf{v}\in\mathcal{V}$, then $T=0$.
\begin{proof}
Set $Q(\mathbf{x})=\innerprod{T\mathbf{x},\mathbf{x}}$.\\
For $\mathbf{u},\mathbf{w}\in\mathcal{V}$, expansion gives the polarization identity for the sesquilinear form $(\mathbf{u},\mathbf{w})\mapsto\innerprod{T\mathbf{u},\mathbf{w}}$:
\begin{align}
\innerprod{T\mathbf{u},\mathbf{w}} =\frac{Q(\mathbf{u}+\mathbf{w})-Q(\mathbf{u}-\mathbf{w})}{4} +i\,\frac{Q(\mathbf{u}+i\mathbf{w})-Q(\mathbf{u}-i\mathbf{w})}{4}. \nonumber
\end{align}
The hypothesis makes every term on the right zero, so $\innerprod{T\mathbf{u},\mathbf{w}}=0$ for all $\mathbf{u},\mathbf{w}$.\\
Taking $\mathbf{w}=T\mathbf{u}$ gives $\norm{T\mathbf{u}}^2=0$, hence $T\mathbf{u}=\mathbf{0}$ for every $\mathbf{u}$.\\
\hlt{Structural Commentary}: Over $\mathbb{C}$, diagonal values of a sesquilinear form determine all mixed values.\\
\hlt{Key Insight}: Polarization turns the one-vector condition into a two-vector condition.
\end{proof}
\end{lemma}

\begin{proposition} \label{prop:complex-quadratic-form-self-adjoint-characterization} \hlt{Self-Adjointness from Real Quadratic Values}\\
Let $\mathcal{V}$ be a finite-dimensional complex inner-product space and let $T\in\LL(\mathcal{V})$.\\
Then $T$ is self-adjoint if and only if $\innerprod{T\mathbf{v},\mathbf{v}}\in\mathbb{R}$ for every $\mathbf{v}\in\mathcal{V}$.
\begin{proof}
For each $\mathbf{v}$, $\innerprod{(T-T^*)\mathbf{v},\mathbf{v}} =\innerprod{T\mathbf{v},\mathbf{v}}-\innerprod{T^*\mathbf{v},\mathbf{v}} =\innerprod{T\mathbf{v},\mathbf{v}}-\overline{\innerprod{T\mathbf{v},\mathbf{v}}}$.\\
If all values $\innerprod{T\mathbf{v},\mathbf{v}}$ are real, this expression is zero for every $\mathbf{v}$, and Lemma \ref{lem:complex-zero-quadratic-form-implies-zero-operator} gives $T-T^*=0$.\\
Conversely, if $T=T^*$, then $\innerprod{T\mathbf{v},\mathbf{v}} =\innerprod{\mathbf{v},T\mathbf{v}} =\overline{\innerprod{T\mathbf{v},\mathbf{v}}}$, so $\innerprod{T\mathbf{v},\mathbf{v}}$ is real.\\
\hlt{Structural Commentary}: The imaginary part of $\innerprod{T\mathbf{v},\mathbf{v}}$ is exactly the diagonal trace of $T-T^*$.\\
\hlt{Key Insight}: Over $\mathbb{C}$, real quadratic values force the operator to equal its adjoint.
\end{proof}
\end{proposition}

\begin{proposition} \label{prop:self-adjoint-zero-quadratic-form-implies-zero} \hlt{Zero Quadratic Form for Self-Adjoint Operators}\\
Let $T$ be a self-adjoint operator on a finite-dimensional inner-product space $\mathcal{V}$.\\
If $\innerprod{T\mathbf{v},\mathbf{v}}=0$ for every $\mathbf{v}\in\mathcal{V}$, then $T=0$.
\begin{proof}
If $\mathbb{F}=\mathbb{C}$, this is Lemma \ref{lem:complex-zero-quadratic-form-implies-zero-operator}. Assume $\mathbb{F}=\mathbb{R}$.\\
For $\mathbf{u},\mathbf{w}\in\mathcal{V}$, self-adjointness gives $\innerprod{T\mathbf{w},\mathbf{u}}=\innerprod{\mathbf{w},T\mathbf{u}}=\innerprod{T\mathbf{u},\mathbf{w}}$, so the real polarization identity gives
\begin{align}
\innerprod{T\mathbf{u},\mathbf{w}}
=\frac{\innerprod{T(\mathbf{u}+\mathbf{w}),\mathbf{u}+\mathbf{w}}-\innerprod{T(\mathbf{u}-\mathbf{w}),\mathbf{u}-\mathbf{w}}}{4}=0.
\nonumber
\end{align}
Taking $\mathbf{w}=T\mathbf{u}$ gives $T\mathbf{u}=\mathbf{0}$ for every $\mathbf{u}$.\\
\hlt{Structural Commentary}: Self-adjointness supplies the missing symmetry needed for real polarization.\\
\hlt{Key Insight}: A self-adjoint operator is determined by its quadratic form.
\end{proof}
\end{proposition}

\begin{definition} \label{def:orthogonal-projection-matrix} \hlt{Orthogonal Projection Matrix}\\
A square matrix $P\in\operatorname{Mat}(n,n,\mathbb{F})$ is an \hlt{orthogonal projection matrix} if $P^2=P$ and $P^*=P$.
\end{definition}

\begin{theorem} \label{thm:orthogonal-projection-self-adjoint-idempotent} \hlt{Orthogonal Projections as Self-Adjoint Idempotents}\\
Let $\mathcal{V}$ be finite-dimensional and let $P\in\LL(\mathcal{V})$.\\
Then $P$ is the orthogonal projection onto $\operatorname{range}(P)$ if and only if $P^2=P$ and $P^*=P$.
\begin{proof}
\begin{enumerate}[label=\roman*.]
\setlength{\itemsep}{0pt}
\item[$(\Rightarrow)$] If $P=P_{\mathcal{U}}$ for $\mathcal{U}=\operatorname{range}(P)$, then Proposition \ref{prop:orthogonal-projection-properties} gives $P^2=P$. For $\mathbf{v}=\mathbf{u}+\mathbf{w}$ and $\mathbf{x}=\mathbf{a}+\mathbf{b}$ with $\mathbf{u},\mathbf{a}\in\mathcal{U}$ and $\mathbf{w},\mathbf{b}\in\mathcal{U}^{\perp}$, $\innerprod{P\mathbf{v},\mathbf{x}}=\innerprod{\mathbf{u},\mathbf{a}+\mathbf{b}}=\innerprod{\mathbf{u},\mathbf{a}} =\innerprod{\mathbf{u}+\mathbf{w},\mathbf{a}}=\innerprod{\mathbf{v},P\mathbf{x}}$, so $P^*=P$.
\item[$(\Leftarrow)$] Suppose $P^2=P$ and $P^*=P$. For each $\mathbf{v}$, $\mathbf{v}=P\mathbf{v}+(\mathbf{v}-P\mathbf{v})$ with $P\mathbf{v}\in\operatorname{range}(P)$ and $P(\mathbf{v}-P\mathbf{v})=\mathbf{0}$.\\
Thus the residual lies in $\ker(P)$. Since $P=P^*$, Proposition \ref{prop:adjoint-kernel-range-orthogonality} gives $\ker(P)=(\operatorname{range}(P))^{\perp}$.\\
Hence every $\mathbf{v}$ has the orthogonal decomposition $\mathbf{v}=P\mathbf{v}+(\mathbf{v}-P\mathbf{v})$ with first part in $\operatorname{range}(P)$, so $P$ is exactly the orthogonal projection onto $\operatorname{range}(P)$.
\end{enumerate}
\hlt{Structural Commentary}: Idempotence splits by range and kernel; self-adjointness makes the split orthogonal.\\
\hlt{Key Insight}: An idempotent is orthogonal exactly when $\ker(P)=(\operatorname{range}(P))^{\perp}$.
\end{proof}
\end{theorem}
